\chapter{Bus Timing Diagrams}
\label{ch:bus-timing}

\section{Bus Cycle Diagrams}

The timing diagrams below show cycle-level logical timing. Each labeled interval represents one CLKI cycle as observed
between rising edges. Waveform rows show physical logic levels; for active-low signals, \texttt{L} means asserted and
\texttt{H} means deasserted. Address, data, \texttt{BAT}, and phase rows are value bands rather than electrical
setup/hold specifications. An empty band means that the row carries no defined value in that cycle. Every cycle begins
with a CLKI rising edge. The pulses in the \texttt{Sample} row mark the rising edges at which the CPU samples
\activelow{BACK}, \activelow{BERR}, \activelow{BPER}, and read data, which are the edges that end the cycles in which
\activelow{BAS} is asserted (Section~\ref{sec:clock-edges}).

Table~\ref{tab:figure-phase-labels} lists the labels used in the phase bands, which name the activity of each cycle.

\begin{table}[H]
	\centering
	\begin{tabularx}{\textwidth}{@{}p{4.2cm}X@{}}
		\toprule
		\textbf{Label}                                                         & \textbf{Meaning}                                                                                                                                  \\
		\midrule
		\texttt{IF high}, \texttt{IF low}                                      & The high-word and low-word instruction-fetch cycles (phases 0 and 1).                                                                             \\
		\texttt{Decode}, \texttt{Execute}, \texttt{Memory}, \texttt{Writeback} & Phases 2, 3, 4 and 5 of an instruction. Section~\ref{sec:execution-phases} defines the phases; no other definition is needed to read the figures. \\
		\texttt{EA}                                                            & The effective-address calculation cycle.                                                                                                          \\
		\texttt{EV high}, \texttt{EV low}                                      & The high-word and low-word cycles of an exception-vector fetch.                                                                                   \\
		\texttt{RV high}, \texttt{RV low}                                      & The high-word and low-word cycles of the reset-vector fetch.                                                                                      \\
		\texttt{Mem}                                                           & The data-access cycle of a load or store.                                                                                                         \\
		\texttt{next}                                                          & The cycle that follows a bus transfer, in which the bus is idle.                                                                                  \\
		\texttt{req}                                                           & In a wait-state figure, the cycle in which the request is first presented.                                                                        \\
		\texttt{wait}                                                          & A cycle in which no response is yet asserted.                                                                                                     \\
		\texttt{done}                                                          & The cycle in which the response completes the transfer.                                                                                           \\
		\texttt{Prev WB}                                                       & The writeback phase of the preceding instruction.                                                                                                 \\
		\texttt{Enter}                                                         & The first cycle of exception entry after the vector is fetched.                                                                                   \\
		\texttt{reset}                                                         & The cycle in which \activelow{RSTI} is first sampled asserted.                                                                                    \\
		\texttt{hold}                                                          & A cycle of the reset-output hold.                                                                                                                 \\
		\texttt{BERR}, \texttt{BPER}                                           & The abort cases.                                                                                                                                  \\
		\bottomrule
	\end{tabularx}
	\caption{Phase Band Labels in the Timing Diagrams}
	\label{tab:figure-phase-labels}
\end{table}

\texttt{BAT} bands name the bus-access type. The diagrams use \texttt{IF} for instruction fetch, \texttt{DR} for data
read, \texttt{DW} for data write, \texttt{EV} for exception vector fetch of vector $n$ at word addresses $2n$ and
$2n+1$, and \texttt{None} for no bus access. \texttt{active} is the type of the access in progress when the abort is
signalled. Address bands show \texttt{PC} and \texttt{PC+1} for the fetch addresses, \texttt{EA} for the effective
address, and \texttt{addr} for the address of the access in progress. Data bands show the \texttt{high} or \texttt{low}
word of a two-word transfer, or \texttt{data} for a single-word transfer.

\begin{figure}[H]
	\centering
	\small
	\begin{tikztimingtable}[sirc timing diagram]
		CLKI        & L 3{0.5H 0.5L}                         \\
		Phase       & D{IF high} D{IF low} D{Decode} D{} \\
		SYNC        & H L L L                       \\
		\activelow{BAS}         & L L H H                       \\
		BRW         & H H H H                       \\
		BAT0--BAT2  & D{IF} D{IF} D{None} D{None}   \\
		A0--A23     & D{PC} D{PC+1} D{} D{}         \\
		D0--D15     & D{high} D{low} Z Z            \\
		\activelow{BACK}        & L L H H                       \\
		\activelow{BERR}/\activelow{BPER}   & H H H H                       \\
		Sample      & 0.95L 0.1H 0.9L 0.1H 1.95L \\
	\end{tikztimingtable}
	\caption{Instruction-Fetch High-Word and Low-Word Reads with No Wait States}
	\label{fig:instruction-fetch-read-timing}
\end{figure}

\begin{figure}[H]
	\centering
	\small
	\begin{tikztimingtable}[sirc timing diagram]
		CLKI        & L 2{0.5H 0.5L}                      \\
		Phase       & D{Mem} D{next} D{}        \\
		SYNC        & L L L                     \\
		\activelow{BAS}         & L H H                     \\
		BRW         & H H H                     \\
		BAT0--BAT2  & D{DR} D{None} D{None}     \\
		A0--A23     & D{EA} D{} D{}             \\
		D0--D15     & D{data} Z Z               \\
		\activelow{BACK}        & L H H                     \\
		\activelow{BERR}/\activelow{BPER}   & H H H                     \\
		Sample      & 0.95L 0.1H 1.95L \\
	\end{tikztimingtable}
	\caption{Data Read Cycle with No Wait State}
	\label{fig:data-read-timing}
\end{figure}

\begin{figure}[H]
	\centering
	\small
	\begin{tikztimingtable}[sirc timing diagram]
		CLKI        & L 2{0.5H 0.5L}                       \\
		Phase       & D{Mem} D{next} D{}         \\
		SYNC        & L L L                      \\
		\activelow{BAS}         & L H H                      \\
		BRW         & L H H                      \\
		BAT0--BAT2  & D{DW} D{None} D{None}      \\
		A0--A23     & D{EA} D{} D{}              \\
		D0--D15     & D{data} Z Z                \\
		\activelow{BACK}        & L H H                      \\
		\activelow{BERR}/\activelow{BPER}   & H H H                      \\
		Sample      & 0.95L 0.1H 1.95L \\
	\end{tikztimingtable}
	\caption{Data Write Cycle with No Wait State}
	\label{fig:data-write-timing}
\end{figure}

\begin{figure}[H]
	\centering
	\small
	\begin{tikztimingtable}[sirc timing diagram]
		CLKI        & L 3{0.5H 0.5L}                                 \\
		Phase       & D{EV high} D{EV low} D{next} D{} \\
		SYNC        & H L L L                              \\
		\activelow{BAS}         & L L H H                              \\
		BRW         & H H H H                              \\
		BAT0--BAT2  & D{EV} D{EV} D{None} D{None}          \\
		A0--A23     & D{$2n$} D{$2n+1$} D{} D{}            \\
		D0--D15     & D{high} D{low} Z Z                   \\
		\activelow{BACK}        & L L H H                              \\
		\activelow{BERR}/\activelow{BPER}   & H H H H                              \\
		Sample      & 0.95L 0.1H 0.9L 0.1H 1.95L \\
	\end{tikztimingtable}
	\caption{Exception-Vector Fetch with High-Word Then Low-Word Reads}
	\label{fig:exception-vector-fetch-timing}
\end{figure}

\begin{figure}[H]
	\centering
	\small
	\begin{tikztimingtable}[sirc timing diagram]
		CLKI        & L 3{0.5H 0.5L}                              \\
		Phase       & D{req} D{wait} D{done} D{next}        \\
		SYNC        & L L L L                           \\
		\activelow{BAS}         & L L L H                           \\
		BRW         & H H H H                           \\
		BAT0--BAT2  & D{DR} D{DR} D{DR} D{None}           \\
		A0--A23     & D{addr} D{addr} D{addr} D{}         \\
		D0--D15     & Z Z D{data} Z                       \\
		\activelow{BACK}        & H H L H                           \\
		\activelow{BERR}/\activelow{BPER}   & H H H H                           \\
		Sample      & 0.95L 0.1H 0.9L 0.1H 0.9L 0.1H 0.95L \\
	\end{tikztimingtable}
	\caption{Wait-State Insertion when No Bus Response Is Asserted}
	\label{fig:bus-wait-state-timing}
\end{figure}

\begin{figure}[H]
	\centering
	\small
	\begin{tikztimingtable}[sirc timing diagram]
		CLKI        & L 3{0.5H 0.5L}                          \\
		Case        & D{BERR} D{} D{BPER} D{}             \\
		\activelow{BAS}         & L H L H                       \\
		BRW         & H H H H                       \\
		\activelow{BACK}        & H H H H                       \\
		\activelow{BERR}        & L H H H                       \\
		\activelow{BPER}        & H H L H                       \\
		BAT0--BAT2  & D{active} D{None} D{active} D{None} \\
		A0--A23     & D{addr} D{} D{addr} D{}             \\
		D0--D15     & Z Z Z Z                       \\
		Sample      & 0.95L 0.1H 1.9L 0.1H 0.95L \\
	\end{tikztimingtable}
	\caption{Bus-Error and Bus-Protection Abort Timing}
	\label{fig:bus-abort-timing}
\end{figure}

\begin{figure}[H]
	\centering
	\small
	\begin{tikztimingtable}[sirc timing diagram]
		CLKI        & L 5{0.5H 0.5L}                                      \\
		Phase       & D{IF high} D{IF low} D{Decode} D{Execute} D{Memory} D{Writeback} \\
		SYNC        & H L L L L L                                \\
		\activelow{BAS}         & L L H H H H                                \\
		BRW         & H H H H H H                                \\
		BAT0--BAT2  & D{IF} D{IF} D{None} D{None} D{None} D{None} \\
		A0--A23     & D{PC} D{PC+1} D{} D{} D{} D{}              \\
		D0--D15     & D{high} D{low} Z Z Z Z                     \\
		\activelow{BACK}        & L L H H H H                                \\
		\activelow{BERR}/\activelow{BPER}   & H H H H H H                                \\
		Sample      & 0.95L 0.1H 0.9L 0.1H 3.95L \\
	\end{tikztimingtable}
	\caption{Complete Register-Only ALU Instruction with No Wait States}
	\label{fig:alu-instruction-cycle-timing}
\end{figure}

\begin{figure}[H]
	\centering
	\small
	\begin{tikztimingtable}[sirc timing diagram]
		CLKI        & L 5{0.5H 0.5L}                                      \\
		Phase       & D{IF high} D{IF low} D{Decode} D{EA} D{Memory} D{Writeback} \\
		SYNC        & H L L L L L                                \\
		\activelow{BAS}         & L L H H L H                                \\
		BRW         & H H H H H H                                \\
		BAT0--BAT2  & D{IF} D{IF} D{None} D{None} D{DR} D{None}  \\
		A0--A23     & D{PC} D{PC+1} D{} D{} D{EA} D{}            \\
		D0--D15     & D{high} D{low} Z Z D{data} Z               \\
		\activelow{BACK}        & L L H H L H                                \\
		\activelow{BERR}/\activelow{BPER}   & H H H H H H                                \\
		Sample      & 0.95L 0.1H 0.9L 0.1H 2.9L 0.1H 0.95L \\
	\end{tikztimingtable}
	\caption{Complete LOAD Instruction with No Wait States}
	\label{fig:load-instruction-cycle-timing}
\end{figure}

\section{Reset and Interrupt Timing Examples}

Figure~\ref{fig:hardware-reset-cycle-timing} shows a hardware reset in which \activelow{RSTI} is asserted for one
clock. The reset-output hold is six CLKI cycles in this case. The cycle in which \activelow{RSTI} is first sampled
asserted counts as the first, and \activelow{RSTO} stays asserted for five further cycles after \activelow{RSTI} is
sampled deasserted. An \activelow{RSTI} assertion of $n$ cycles therefore gives $n + 5$ cycles of \activelow{RSTO}. A
software reset and a reset-vector fetch retry each give a six-cycle hold. The reset-vector fetch does not begin until
\activelow{RSTI} is deasserted and the reset-output hold has completed.

\begin{figure}[H]
	\centering
	\small
	\begin{tikztimingtable}[sirc timing diagram]
		CLKI        & L 7{0.5H 0.5L}                                         \\
		Phase       & D{reset} D{hold} D{hold} D{hold} D{hold} D{hold} D{RV high} D{RV low} \\
		\activelow{RSTI}        & L H H H H H H H                               \\
		\activelow{RSTO}        & L L L L L L H H                               \\
		\activelow{BAS}         & H H H H H H L L                               \\
		BRW         & H H H H H H H H                               \\
		BAT0--BAT2  & D{None} D{None} D{None} D{None} D{None} D{None} D{EV} D{EV} \\
		A0--A23     & D{} D{} D{} D{} D{} D{} D{0} D{1}     \\
		D0--D15     & Z Z Z Z Z Z D{high} D{low}                    \\
		\activelow{BACK}        & H H H H H H L L                               \\
		\activelow{BERR}/\activelow{BPER}   & H H H H H H H H                               \\
		Sample      & 6.95L 0.1H 0.9L 0.05H \\
	\end{tikztimingtable}
	\caption{Hardware Reset with One-Cycle \activelow{RSTI} Assertion and Reset-Vector Fetch}
	\label{fig:hardware-reset-cycle-timing}
\end{figure}

Figure~\ref{fig:irq-sampling-cycle-timing} shows an enabled maskable interrupt asserted before an instruction boundary.
The CPU samples the interrupt lines on every clock cycle, so the request is latched in the cycle it is first seen, but
it is dispatched only at the next phase 0 at which no coprocessor call or command is pending or active. A request first
sampled at the rising edge that starts phase 0 is dispatched in that phase 0. The following exception-unit dispatch
fetches the selected hardware-interrupt vector instead of fetching the next processing-unit instruction. A request
latched partway through an instruction waits for the next such phase 0. The figure holds
\activelow{IRQ1}--\activelow{IRQ4} asserted throughout because the requesting device must hold the line asserted until
its handler acknowledges it.

\begin{figure}[H]
	\centering
	\small
	\begin{tikztimingtable}[sirc timing diagram]
		CLKI        & L 4{0.5H 0.5L}                                  \\
		Phase       & D{Prev WB} D{EV high} D{EV low} D{Decode} D{Enter} \\
		\activelow{IRQ1}--\activelow{IRQ4}  & L L L L L                              \\
		SYNC        & L H L L L                              \\
		\activelow{BAS}         & H L L H H                              \\
		BRW         & H H H H H                              \\
		BAT0--BAT2  & D{None} D{EV} D{EV} D{None} D{None}    \\
		A0--A23     & D{} D{$2n$} D{$2n+1$} D{} D{}          \\
		D0--D15     & Z D{high} D{low} Z Z                   \\
		\activelow{BACK}        & H L L H H                              \\
		\activelow{BERR}/\activelow{BPER}   & H H H H H                              \\
		Sample      & 1.95L 0.1H 0.9L 0.1H 1.95L \\
	\end{tikztimingtable}
	\caption{Enabled \activelow{IRQ} Latched Before an Instruction Boundary and Dispatched with No Wait States}
	\label{fig:irq-sampling-cycle-timing}
\end{figure}
